% Verbatim copy of /Users/alingalatan/workspaces/grothe/handoff/manuscript/legacy/lower-bound/sharp-influence.tex (file dated 2026-10-25). Only this line was added.
\subsubsection{The sharp global influence coefficient: reconciliation and proof}
\label{sec:sharp-influence-audit}
\begin{theorem}[Sharp global influence bound, \st{P,C}]\label{thm:sharp-influence-audit}
For a finite-dimensional pure complex cubic $q$, let
\[
 e=\E|q|^2,\quad \kappa=\E|q|^4-2e^2,\quad
 I_q=\E[dd^*]+\E[\bar f f^{\mathsf T}],\quad
 d_j=\partial_{X_j}q,\quad f_j=\partial_{\bar X_j}q.
\]
Thus $I_q$ is the matrix $K$ of (SB.6a), not its restriction to the
shared-driver span. Then
\[
 \boxed{\kappa\le\frac{14}{3}\operatorname{tr}(I_q^2)
                 \le14e\|I_q\|_{\rm op}.}
\]
Both constants are sharp. No sixth-moment assumption is used.
\end{theorem}
\begin{proof}
Write $q=\sum t_{abc}{:}X_aX_b\bar X_c{:}$, $t_{abc}=t_{bac}$,
so $e=2\|t\|^2$. Define
\[
 N_{ab}=\sum_{j,k}t_{ajk}\bar t_{bjk},\qquad
 P_{ab}=\sum_{j,k}t_{jka}\bar t_{jkb},\qquad
 B_{be}=\sum_{a,d}t_{abd}t_{dea},
\]
\[
 U_{abj;c}=\sum_i t_{abi}t_{jic},\qquad
 X=\operatorname{Sym}_{a,b,j}U,\qquad z=\operatorname{tr}(NP^{\mathsf T}).
\]
The derivative normalization gives $I_q=4N+2P^{\mathsf T}$ and
$\operatorname{tr}I_q=3e$. Wick expansion gives
\[
 \kappa=32\|N\|_{\rm HS}^2+8\|P\|_{\rm HS}^2
                    +96\|X\|^2+32\|B\|_{\rm HS}^2.             \tag{SI.1}
\]
The exact checker enumerates all 176 pairings of the six holomorphic
and six antiholomorphic legs of $q,q,\bar q,\bar q$, excluding pairs
within one Wick factor. Eight disconnected pairings contribute $2e^2$.
For the remainder it identifies contraction graphs under exchanges of
identical factors, exchanges of the first two tensor slots, and renaming
of summed indices. The six resulting graph types agree coefficientwise
with (SI.1), including the nine terms of $\|\operatorname{Sym}U\|^2$.
This is an identity of finite sums for every index set, rather than a
numerical check in finitely many dimensions.

Expanding the squared Gram norm now gives the exact slack identity
\begin{align*}
 \frac{14}{3}\|I_q\|_{\rm HS}^2-\kappa
 ={}&32(\|N\|_{\rm HS}^2-\|B\|_{\rm HS}^2)
       +96(z-\|X\|^2)\\
    &+\frac{32}{3}\|N-P^{\mathsf T}\|_{\rm HS}^2.                \tag{SI.2}
\end{align*}
All three terms are nonnegative. Index rearrangement gives
$\|U\|^2=z$, and symmetrization is an orthogonal projection.
For the first term, let the columns of a matrix $R$ be the vectorized
slices $(t_{abd})_{a,d}$ indexed by $b$, and let $J$ transpose a slice.
Then $B=R^{\mathsf T}JR$ and the Gram matrix $R^*R$ is $N^{\mathsf T}$.
The Schatten--H\"older inequality gives
\[
 \|B\|_{
\rm HS}\le\|R^{\mathsf T}\|_4\|JR\|_4
 =\|R\|_4^2=\|N\|_{\rm HS}.
\]
Finally $I_q\succeq0$ and $\operatorname{tr}I_q=3e$ imply
$\operatorname{tr}(I_q^2)\le3e\|I_q\|_{\rm op}$.
For $q=N^{-1/2}\sum_{j=1}^NH(X_j)$,
\[
 e=1,\quad\kappa=42/N,\quad I_q=(3/N)\operatorname{Id}_N,
 \quad\operatorname{tr}(I_q^2)=9/N.
\]
Both inequalities are equalities. Unequal independent Wick sums still
saturate the first inequality; equality in the second requires equal
nonzero squared amplitudes.
\end{proof}

\paragraph{The fourth-moment cap is a direct consequence.}
Since $h\le\operatorname{tr}I_q=3e$, the same theorem gives
$\E|q|^4\le44e^2$. For $e=1$ this is stronger than the cap
$2+(31/4)^2=993/16$ assumed by \code{Bridge.gaussian\_moments}.
For the intended $H_{2,1}$ residual, this supplies the mathematical cap.
The Lean bridge is stated more broadly for two real third-chaos
polynomials: its signature alone does not encode bidegree $(2,1)$.
A port must retain that extra structure rather than apply the bound to
an arbitrary pair of real cubic chaoses.

\paragraph{What is new in this reconciliation.}
The value $14$ was already recorded in \S\ref{sec:wct-external}; the later
real-chaos route in (SB.6) retained the weaker $50/3$. The proof and exact
contraction check above give one canonical statement with consistent
normalizations. Keep the real-matrix bound as a separate option: it can
be stronger for particular derivative covariance matrices. Its numerical
right-hand side need not itself be bounded by $14eh$.

\begin{corollary}[A larger diffuse-tensor region]\label{cor:sharp-influence-region}
For the full coupled A/B family (CT.1), with arbitrary phases and
amplitudes, the assumptions
\[
 g\ge9/40,\qquad e\le1/25,\qquad h\le e/10
\]
imply
$\Lambda(G)<S-1/40000-E_{AB}/2000$.
Also $g\ge1/5$, $e\le1/25$, $h\le e/150$ imply (CT.4).
\end{corollary}
\begin{proof}
The original numerical certificate behind Theorem~\ref{thm:ct-influence} treats $h$ and $\kappa$ separately.
For its exact branch, $14eh\le7e^2/5$; for its weak branch,
$14eh\le7e^2/75<e^2/10$. These verify every hypothesis without
rerunning or enlarging the earlier numerical certificate.
\end{proof}
The exact branch extends the previously stated influence-only cutoff
$9e/100$ to $e/10$, with the displayed (smaller) guaranteed deficit.
The weak branch inherits its accepted remote comparison. By themselves these A/B arguments do not include arbitrary even sectors.
Corollary~\ref{cor:quadratic-all-sector-refinement} supplies their simultaneous
even/diagonal extension with explicit penalties. A complete row envelope
is still missing.
